\chapter{Methodology}
\label{ch:methodology}

This chapter presents the evaluation framework developed in this thesis.
Section~\ref{sec:eval-setup} defines the evaluation objective and the unit of
analysis. Sections~\ref{sec:dataset}--\ref{sec:fae-methods} describe the
dataset, predictive models, and feature attribution methods that constitute
the experimental infrastructure. Section~\ref{sec:metrics} specifies the
twelve evaluation metrics organized across six Quantus categories.
Sections~\ref{sec:normalization}--\ref{sec:aggregation} detail the score
normalization and aggregation procedures that form the core methodological
contribution of this work (Contributions~1--3, per
Section~\ref{sec:celePracy}). Finally, Section~\ref{sec:pipeline} describes
the end-to-end evaluation pipeline and its software implementation.

%% ========================================================================
\section{Evaluation Setup}
\label{sec:eval-setup}

The objective of this evaluation framework is to rank Feature Attribution
Explanation (FAE) methods by an aggregated effectiveness index
$E(\Phi)$ that summarizes performance across six metric categories:
Faithfulness, Robustness, Localization, Complexity, Randomization, and
Axiomatic.

The atomic unit of evaluation is a \emph{per-(model, image, FAE method)
attribution triple}. For a classifier $f$, an input image $x$, and an
explanation method $\Phi$, the triple $(f, x, \Phi)$ yields an attribution
map $\Phi(f, x) \in \mathbb{R}^{C \times H \times W}$, where $C$, $H$, $W$
denote the channel, height, and width dimensions of the input.

For each triple, the \textbf{target class} used for attribution is the
model's top-1 prediction $\hat{y} = \operatorname{argmax}_c f(x)_c$, not
the ground-truth label. This choice is deliberate: evaluation metrics assess
explanation fidelity to \emph{the model's decision}, not to a human-supplied
label. An explanation that faithfully reflects the features driving the
model's prediction is informative regardless of whether that prediction is
correct \cite{hedstrom2023quantus}.

%% ========================================================================
\section{Dataset}
\label{sec:dataset}

The experiments use the ISIC 2017 Challenge dataset (Task~1 --- Lesion
Segmentation), publicly available under a CC-BY-NC license from
\texttt{challenge.isic-archive.com}. The dataset comprises dermoscopic
images distributed across three diagnostic categories: melanoma, nevus, and
seborrheic keratosis.

\begin{itemize}
    \item \textbf{Split sizes:} 2{,}000 training images, 150 validation
    images, 600 test images.
    \item \textbf{Segmentation masks:} Each image is accompanied by an
    expert-annotated binary mask delineating the lesion boundary. These masks
    enable the computation of Localization metrics (Pointing Game, Relevance
    Mass Accuracy) without additional labeling effort---a key advantage over
    datasets that provide only image-level labels.
    \item \textbf{Preprocessing:} All images are resized to $224 \times 224$
    pixels and normalized using ImageNet channel-wise mean and standard
    deviation ($\mu = [0.485, 0.456, 0.406]$, $\sigma = [0.229, 0.224,
    0.225]$).
\end{itemize}

The vertical-slice validation (Section~\ref{sec:pipeline}) uses a subset of
12 images from the test split. The full experiment scales to all 600 test
images.
\studentinline{TODO: Update image counts and confirm split sizes after full experiment run.}

%% ========================================================================
\section{Predictive Models}
\label{sec:models}

Two convolutional neural network architectures serve as the predictors whose
decisions are explained. They are not a contribution of this thesis; they
exist solely to produce the attributions that the evaluation framework then
scores.

\subsection{ResNet-18}

ResNet-18 is a residual network with 18 layers organized in four residual
blocks with skip connections \cite{he2016resnet}.
\studentinline{VERIFY citation: He et al.\ 2016 --- ``Deep Residual Learning for Image Recognition.''} The final
fully connected layer is replaced with $\text{Linear}(512, 3)$ to match the
three-class ISIC 2017 classification task.

\subsection{SqueezeNet 1.1}

SqueezeNet 1.1 is a compact architecture that achieves AlexNet-level accuracy
with $\sim$50$\times$ fewer parameters, using Fire modules that combine
$1 \times 1$ squeeze convolutions with mixed $1 \times 1$ and $3 \times 3$
expand convolutions \cite{iandola2016squeeze}.
\studentinline{VERIFY citation: Iandola et al.\ 2016 --- ``SqueezeNet: AlexNet-Level Accuracy with 50x Fewer Parameters.''} The final classifier layer is replaced with
$\text{Conv2d}(512, 3, 1, 1)$ for three-class output.

\subsection{Training Protocol}

Both models are trained with the following configuration:
\begin{itemize}
    \item \textbf{Epochs:} 25
    \item \textbf{Optimizer:} SGD with momentum ($\text{lr} = 0.001$,
    $\text{momentum} = 0.9$)
    \item \textbf{Scheduler:} StepLR (step size 7, $\gamma = 0.1$)
    \item \textbf{Loss:} Class-weighted CrossEntropyLoss to address class
    imbalance in the ISIC 2017 dataset
    \item \textbf{Random seed:} 42 (set for \texttt{torch}, \texttt{numpy},
    and \texttt{random})
\end{itemize}

Best validation accuracies: ResNet-18 achieves 75.3\%, SqueezeNet achieves
74.0\%. Trained weights are saved to disk with SHA-256 checksums documented
in \texttt{weights/README.md}.

%% ========================================================================
\section{Feature Attribution Methods}
\label{sec:fae-methods}

Seven attribution methods are evaluated, spanning gradient-based (white-box)
and perturbation-based (black-box) families. All methods are implemented
using the Captum library, ensuring consistent interfaces and reproducible
gradient computation.

\subsection{Gradient-based Methods}

\subsubsection{Saliency (Vanilla Gradient)}
Computes the absolute gradient of the class score with respect to the input
image (Section~\ref{ch:sota}). Computationally efficient but prone to noisy,
non-class-discriminative maps due to shattered gradients.

\subsubsection{Integrated Gradients}
Accumulates gradients along a linear path from a black baseline ($x' = 0$)
to the actual input $x$ using $n\_\text{steps} = 50$ interpolation points
\cite{sundararajan2017axiomatic}. Satisfies the Completeness axiom by
construction.

\subsubsection{Grad-CAM}
Gradient-weighted Class Activation Mapping applied to the last convolutional
layer of the network \cite{selvaraju2017grad}. Produces a coarse
localization map at the resolution of the final feature map, which is then
upsampled to input resolution via bilinear interpolation.

\subsubsection{Guided Backpropagation}
Propagates gradients through only the positive activations at each ReLU
layer, producing sharper maps than vanilla gradients but at the risk of
acting as an edge detector independent of model parameters
\cite{springenberg2015}.
\studentinline{VERIFY citation: Springenberg et al.\ 2015 --- ``Striving for Simplicity: The All Convolutional Net.''}

\subsubsection{DeepLift}
Computes attribution by comparing each neuron's activation to a reference
activation produced by a black baseline input ($x' = 0$). Decomposes the
output difference $f(x) - f(x')$ into individual feature contributions
\cite{shrikumar2017deeplift}.
\studentinline{VERIFY citation: Shrikumar et al.\ 2017 --- ``Learning Important Features Through Propagating Activation Differences.''}

\subsubsection{Layer-wise Relevance Propagation (LRP)}
Redistributes the model output backward through the network according to
propagation rules that preserve the total relevance at each layer
\cite{bach2015lrp}.
\studentinline{VERIFY citation: Bach et al.\ 2015 --- ``On Pixel-Wise Explanations for Non-Linear Classifier Decisions by Layer-Wise Relevance Propagation.''}

\subsection{Perturbation-based Methods}

\subsubsection{Occlusion}
A sliding window of size $(3, 15, 15)$ with stride $(3, 8, 8)$ is passed
over the input, replacing occluded regions with a black baseline. The change
in model output for each position yields a sensitivity map
\cite{zeiler2014}.
\studentinline{VERIFY citation: Zeiler \& Fergus 2014 --- ``Visualizing and Understanding Convolutional Networks.''}

\medskip
\noindent
LIME and KernelSHAP are in scope for the full experiment but are excluded
from the vertical slice due to their substantially higher runtime cost.
\studentinline{TODO: Update this section if LIME/KernelSHAP are added in the full experiment.}

\subsection{Attribution Caching}
\label{subsec:caching}

Attributions are cached to disk, keyed by a tuple of (model architecture,
FAE method, image identifier, target class, hyperparameters hash). This
avoids redundant recomputation when the same attribution is needed by
multiple metrics or across repeated pipeline runs.

The cache is \textbf{bypassed} for inner re-computations performed by
robustness and randomization metrics---specifically Max-Sensitivity,
Avg-Sensitivity, and MPRT---because these metrics intentionally perturb
inputs or model weights and therefore produce legitimately different
attributions on each invocation. Caching such perturbed attributions would
violate the metric's operational semantics.

%% ========================================================================
\section{Metric Computation}
\label{sec:metrics}

Twelve metrics across six Quantus categories are computed for each
attribution triple. Table~\ref{tab:metrics} summarizes their directionality
and key hyperparameters.

\begin{table}[H]
\centering
\caption{Evaluation metrics: category, direction, and hyperparameters.}
\begin{tabularx}{\linewidth}{llcl}
\toprule
Category & Metric & Direction & Key hyperparameters \\
\midrule
Faithfulness & FaithfulnessCorrelation & $\uparrow$ & nr\_runs=100, subset\_size=224 \\
Faithfulness & PixelFlipping           & $\uparrow$ & features\_in\_step=224 (AUC) \\
Robustness   & MaxSensitivity          & $\downarrow$ & nr\_samples=10, lower\_bound=0.2 \\
Robustness   & AvgSensitivity          & $\downarrow$ & nr\_samples=10, lower\_bound=0.2 \\
Localization & RelevanceMassAccuracy   & $\uparrow$ & requires mask \\
Localization & PointingGame            & $\uparrow$ & requires mask \\
Complexity   & Sparseness              & $\uparrow$ & Gini-like coefficient \\
Complexity   & Complexity              & $\downarrow$ & Shannon entropy \\
Randomization & ModelParameterRandomisation & $\downarrow$ & layer\_order=top\_down \\
Randomization & RandomLogit            & $\downarrow$ & num\_classes=3 \\
Axiomatic    & Completeness            & $\downarrow$ & black baseline; IG/DeepLift/LRP only \\
Axiomatic    & NonSensitivity          & $\downarrow$ & downsampled to 56$\times$56; excluded by default \\
\bottomrule
\end{tabularx}
\label{tab:metrics}
\end{table}

\paragraph{Faithfulness.}
FaithfulnessCorrelation and PixelFlipping both quantify how well the
attribution map identifies features that the model relies on for its
prediction. FaithfulnessCorrelation uses random subset masking with 100
runs to estimate the Pearson correlation between attribution magnitude and
output change. PixelFlipping removes features in descending order of
attributed importance and aggregates the output degradation as an AUC
score. Both metrics are higher-is-better \cite{hedstrom2023quantus}.

\paragraph{Robustness.}
MaxSensitivity and AvgSensitivity evaluate explanation stability under
small Gaussian perturbations of the input. Both use 10 perturbation samples
with a perturbation magnitude lower bound of 0.2. MaxSensitivity reports
the worst-case sensitivity, while AvgSensitivity reports the average.
Lower is better for both \cite{hedstrom2023quantus}.

\paragraph{Localization.}
RelevanceMassAccuracy and PointingGame evaluate spatial agreement between
the attribution map and the expert-annotated lesion mask. Both are
higher-is-better and require segmentation masks at inference time
\cite{hedstrom2023quantus}.

\paragraph{Complexity.}
Sparseness (Gini-like, higher-is-better) and Complexity (Shannon entropy,
lower-is-better) assess the conciseness of the attribution distribution.
These metrics are model-agnostic and computed directly from the attribution
map \cite{hedstrom2023quantus}.

\paragraph{Randomization.}
ModelParameterRandomisation (MPRT) progressively randomizes model weights
top-down and measures correlation decay. RandomLogit replaces the target
class and measures explanation change. Lower correlation is better for both,
indicating that the explanation genuinely depends on learned parameters and
the specified target class \cite{hedstrom2023quantus}.

\paragraph{Axiomatic.}
Completeness measures the absolute deviation from the completeness axiom
($|\sum_i E_i - \Delta f| \approx 0$). It returns NaN for methods outside
\{Integrated Gradients, DeepLift, LRP\}; these NaN values are excluded
from aggregation rather than imputed. NonSensitivity measures attribution
mass on provably non-contributory features. Due to its $O(d)$ forward-pass
cost, the input is downsampled to $56 \times 56$, and the metric is
excluded from vertical-slice runs by default \cite{hedstrom2023quantus}.

%% ========================================================================
\section{Score Normalization}
\label{sec:normalization}

The twelve metrics operate on heterogeneous scales (e.g.,
FaithfulnessCorrelation $\in [-1, 1]$, MaxSensitivity $\in [0, \infty)$,
Sparseness $\in [0, 1]$) and have heterogeneous directionalities (some
higher-is-better, some lower-is-better). Aggregation requires
rectification and rescaling.

\subsection{Direction Rectification}

Each metric $k$ is assigned a direction sign $s_k \in \{+1, -1\}$, where
$s_k = +1$ for higher-is-better metrics and $s_k = -1$ for
lower-is-better metrics. The rectified score is:
\begin{equation}
    \hat{M}_k(\Phi) = s_k \cdot M_k(\Phi),
\end{equation}
so that a higher rectified score always indicates better performance.

\subsection{Second Moment Scaling}

The default scale-rectifying transformation is Second Moment Scaling,
adopted from the NormEnsembleXAI framework (Decision~D1)
\cite{hryniewska2024normensemblexai}:
\begin{equation}
    \tilde{M}_k(\Phi) = \frac{\hat{M}_k(\Phi)}{\sqrt{\mathbb{E}\bigl[\hat{M}_k^2\bigr]}},
\end{equation}
where the expectation is computed within each (model, metric) group
independently. This transformation does not require knowledge of global
extremes, is robust to outliers (unlike min--max normalization), and
preserves the sign structure of the rectified scores.

Two alternative normalizations are implemented for ablation studies:
\begin{itemize}
    \item \textbf{$z$-score normalization:}
    $\psi_k(x) = (x - \mu_k) / \sigma_k$, centering on the group mean with
    unit variance.
    \item \textbf{Min--max normalization:}
    $\psi_k(x) = (x - \min_k) / (\max_k - \min_k)$, mapping scores to
    $[0, 1]$. Sensitive to outliers and requires knowing the global extremes.
\end{itemize}

%% ========================================================================
\section{Score Aggregation}
\label{sec:aggregation}

Once normalized, the per-metric scores are aggregated into a single
effectiveness index $E(\Phi)$ per (model, image, FAE method). Three
aggregation operators are implemented (Decision~D2):

\begin{enumerate}
    \item \textbf{Weighted mean (default, recommended):}
    \begin{equation}
        E(\Phi) = \sum_{k=1}^{K} w_k \, \tilde{M}_k(\Phi),
    \end{equation}
    where $w_k \geq 0$ and $\sum_k w_k = 1$.

    \item \textbf{Minimum (conservative):}
    \begin{equation}
        E(\Phi) = \min_{k=1,\dots,K} \tilde{M}_k(\Phi).
    \end{equation}
    Evaluates a method by its worst dimension. In safety-critical domains, a
    method that is faithful but unstable is ineligible regardless of how high
    its faithfulness score is.

    \item \textbf{Geometric mean:}
    \begin{equation}
        E(\Phi) = \exp\!\left(\sum_{k=1}^{K} w_k \ln\bigl(\tilde{M}_k(\Phi) + \varepsilon\bigr)\right),
    \end{equation}
    where $\varepsilon > 0$ is a small constant to avoid $\ln(0)$. Penalizes
    methods with a near-zero score on any single axis.
\end{enumerate}

The default weight vector uses uniform weights: $w_k = 1/K$ for $K$
retained metrics after the redundancy and meta-evaluation filtering
described in Chapter~\ref{ch:sota}.

\subsection{Autoweighted Extension}
\label{sec:weighting}

The Autoweighted mechanism, inspired by \citet{hryniewska2024normensemblexai},
derives weights from each metric's ensemble score $\mathrm{ES}_k$:
\begin{equation}
    w_k = \frac{\mathrm{ES}_k}{\sum_{j} \mathrm{ES}_j}.
\end{equation}
In this thesis, $\mathrm{ES}_k$ denotes the MetaQuantus-discounted ensemble
score: the raw ensemble agreement of metric $k$ is attenuated by its
reliability as assessed by the Noise Resilience and Adversarial Reactivity
tests of \citet{hedstrom2024metaquantus}. This ensures that unreliable
metrics receive lower weight in the final aggregation, even if their raw
ensemble scores are high.

\subsection{MetaQuantus-discounted Weighting (Contribution~2)}
\label{subsec:mq-discount}

The novel weighting scheme of this thesis (Decision~D3) combines two
signals that are conventionally treated in isolation: how strongly a metric
agrees with the cross-metric consensus, and how reliable that metric is when
probed by meta-evaluation. The derivation proceeds in three steps.

\paragraph{Step 1 --- Raw agreement weight.}
For each retained metric $k \in M^{*}$, let $\mathrm{ES}_k$ be the
\emph{ensemble agreement} of metric $k$, defined as the mean absolute
rank correlation between the per-method ranking induced by metric $k$ alone
and the consensus ranking induced by the (uniform) aggregate of all metrics
in $M^{*}$:
\begin{equation}
    \mathrm{ES}_k = \frac{1}{|M^{*}| - 1}
        \sum_{j \in M^{*},\, j \neq k}
        \bigl| \rho\bigl(\mathbf{r}_k,\, \mathbf{r}_j\bigr) \bigr|,
\end{equation}
where $\mathbf{r}_k$ is the vector of method ranks under metric $k$ and
$\rho$ is the Spearman rank correlation. A metric whose ranking aligns with
the rest of the panel receives a high $\mathrm{ES}_k$; an idiosyncratic
metric receives a low one.

\paragraph{Step 2 --- Reliability discount.}
Each metric is assigned a reliability coefficient $r_k \in [0, 1]$ obtained
from the meta-evaluation stage (Section~\ref{subsec:meta-evaluation} of
Chapter~\ref{ch:sota}). MetaQuantus produces two scores per metric: a Noise
Resilience score $\mathrm{NR}_k$ and an Adversarial Reactivity score
$\mathrm{AR}_k$. Their combination
\begin{equation}
    r_k = \tfrac{1}{2}\bigl(\mathrm{NR}_k + \mathrm{AR}_k\bigr)
\end{equation}
is the per-metric \texttt{combined\_reliability} reported in
\texttt{results/meta\_evaluation\_reliability.csv}, averaged over the
(model, FAE) cells in which the metric is defined. Where only one of the two
sub-scores is available, $r_k$ falls back to that sub-score.

\paragraph{Step 3 --- Discounted normalisation.}
The agreement weight is attenuated by the reliability coefficient and
renormalised to a convex combination:
\begin{equation}
    w_k = \frac{r_k \, \mathrm{ES}_k}{\sum_{j \in M^{*}} r_j \, \mathrm{ES}_j},
    \qquad w_k \geq 0, \quad \sum_{k \in M^{*}} w_k = 1.
    \label{eq:mq-discount}
\end{equation}
Equation~\eqref{eq:mq-discount} is the MetaQuantus-discounted Autoweighted
operator. It reduces to uniform weighting when all $r_k \mathrm{ES}_k$ are
equal, to the plain Autoweighted scheme of
\citet{hryniewska2024normensemblexai} when every metric is perfectly
reliable ($r_k \equiv 1$), and it drives $w_k \to 0$ for any metric that
either disagrees with the consensus or fails meta-evaluation. The three
weighting schemes---uniform, Autoweighted, and MQ-discounted---are reported
side by side in Chapter~\ref{ch:experiments} so that their effect on the
final method ranking can be assessed directly.

%% ========================================================================
\section{The Evaluation Pipeline}
\label{sec:pipeline}

The end-to-end evaluation pipeline integrates all components described in
this chapter into a single reproducible data flow:

\begin{enumerate}
    \item \textbf{Input:} ISIC 2017 test images and associated binary lesion
    masks.
    \item \textbf{Model inference:} For each image $x$, the classifier $f$
    produces a prediction $\hat{y} = \operatorname{argmax}_c f(x)_c$, which
    serves as the target class for attribution.
    \item \textbf{Attribution generation:} Each FAE method $\Phi$ computes an
    attribution map $\Phi(f, x) \in \mathbb{R}^{3 \times 224 \times 224}$.
    Results are cached to disk (Section~\ref{subsec:caching}).
    \item \textbf{Metric computation:} The twelve metrics from
    Table~\ref{tab:metrics} are applied to each attribution map, producing a
    vector of raw scores per (model, image, FAE method).
    \item \textbf{Normalization:} Direction rectification and Second Moment
    Scaling (Section~\ref{sec:normalization}) transform raw scores into a
    comparable scale.
    \item \textbf{Aggregation:} The normalized scores are combined into a
    single effectiveness index $E(\Phi)$ via the chosen aggregation operator
    (Section~\ref{sec:aggregation}).
    \item \textbf{Ranking:} FAE methods are ranked per-model by their mean
    $E(\Phi)$ across test images.
\end{enumerate}

Figure~\ref{fig:pipeline} summarises this flow. The diagram makes explicit
the two contribution boundaries: the meta-evaluation block
(redundancy~+~meta-validation) yields the reduced metric set $M^{*}$
(Contribution~1), and the aggregation block (normalisation, weighting,
combination) yields the effectiveness index $E(\Phi)$ (Contribution~2),
which the statistical-comparison stage (Contributions~3--4) consumes.

\begin{figure}[H]
    \centering
    \begin{tikzpicture}[
        node distance=5mm and 0mm,
        box/.style={draw, rounded corners, align=center, inner sep=4pt,
                    text width=0.74\linewidth, font=\small},
        contrib/.style={draw, rounded corners, align=center, inner sep=4pt,
                    text width=0.74\linewidth, font=\small, fill=black!5},
        arr/.style={-{Latex[length=2mm]}, thick}
    ]
        \node[box]     (data)    {ISIC 2017 images $+$ expert lesion masks};
        \node[box,below=of data]    (train)   {Base classifier training\\(ResNet-18, SqueezeNet 1.1)};
        \node[box,below=of train]   (attr)    {Attribution generation (Captum, 7 FAE methods)\\$\Phi(f,x)$, cached to disk};
        \node[box,below=of attr]    (metric)  {Metric computation (Quantus, 12 metrics / 6 categories)};
        \node[contrib,below=of metric] (meta) {\textbf{Meta-evaluation (Contribution~1)}\\redundancy pruning $+$ MetaQuantus NR/AR\\$\Rightarrow$ reduced validated set $M^{*}$};
        \node[contrib,below=of meta]   (agg)  {\textbf{Aggregation (Contribution~2)}\\normalise $\to$ rectify $\to$ weight (MQ-discount) $\to$ combine\\$\Rightarrow$ effectiveness index $E(\Phi)$};
        \node[contrib,below=of agg]    (stat) {\textbf{Statistical comparison (Contributions~3--4)}\\Friedman $+$ Nemenyi, Wilcoxon;\\individual vs.\ ensembled attributions};

        \draw[arr] (data)   -- (train);
        \draw[arr] (train)  -- (attr);
        \draw[arr] (attr)   -- (metric);
        \draw[arr] (metric) -- (meta);
        \draw[arr] (meta)   -- (agg);
        \draw[arr] (agg)    -- (stat);
    \end{tikzpicture}
    \caption{End-to-end evaluation pipeline. White boxes are experimental
    infrastructure; shaded boxes are the methodological contributions of this
    thesis.}
    \label{fig:pipeline}
\end{figure}

\subsection{Implementation}

The pipeline entry point is \texttt{experiments/vertical\_slice.py}. The
output is a CSV file with columns: \texttt{model}, \texttt{image\_id},
\texttt{fae\_method}, \texttt{metric}, \texttt{score}. For the 12-image
vertical slice, the expected output is $2 \times 7 \times 12 \times 12 =
2{,}016$ rows (2 models $\times$ 7~FAE methods $\times$ 12~images $\times$
12~metrics).

\studentinline{TODO: Update row count and CSV column description after full experiment run.}

The software architecture follows the modular layout described in
Section~\ref{sec:celePracy}:
\begin{itemize}
    \item \texttt{src/data/} --- ISIC 2017 dataset class and download script.
    \item \texttt{src/models/} --- Classifier definitions and training loop.
    \item \texttt{src/attributions/} --- Attribution generation, disk cache,
    NormEnsembleXAI wrapper.
    \item \texttt{src/metrics/} --- Quantus wrapper, custom metrics,
    redundancy analysis.
    \item \texttt{src/meta\_evaluation/} --- MetaQuantus wrapper and metric
    selection.
    \item \texttt{src/aggregation/} --- Normalization, aggregation operators,
    weighting schemes.
    \item \texttt{src/comparison/} --- Wilcoxon, Friedman, and Nemenyi
    statistical tests.
    \item \texttt{src/pipeline.py} --- End-to-end orchestrator.
\end{itemize}

All experiments are configured via YAML files under \texttt{configs/}, with
one configuration per experiment variant. Random seeds, model paths, metric
hyperparameters, and normalization choices are specified declaratively to
ensure full reproducibility.
